\chapter{Strutture dati e tipi di dato astratti}
\label{chap:strutture-dati-adt}

Gli algoritmi di ricerca e ordinamento del \cref{chap:ricerca-ordinamento}
operano su collezioni di dati. La scelta di come organizzarle dipende
sia dalle operazioni richieste dal problema, sia dal costo della loro
esecuzione, studiato nel \cref{chap:complessita}. Si distinguono pertanto
il comportamento del contenitore e la sua rappresentazione interna.\footnote{Materiale
di riferimento: G. Baruzzo, \emph{Strutture dati e tipi di dati astratti},
dispense del corso di Fondamenti di Informatica.}

\section{Tipi di dato astratti e strutture dati}
\label{sec:fi-adt-strutture}

\begin{definition}[Tipo di dato astratto]
\label{def:fi-adt}
Un \emph{tipo di dato astratto} (\emph{Abstract Data Type}, ADT)
specifica un insieme di valori e le operazioni applicabili a essi,
definendone il comportamento, i risultati, le precondizioni e i casi
di errore, senza stabilire come i dati siano rappresentati né come
le operazioni siano implementate.
\end{definition}

\begin{definition}[Struttura dati]
\label{def:fi-struttura-dati}
Una \emph{struttura dati} stabilisce come gli elementi e le
informazioni necessarie a organizzarli siano rappresentati e
memorizzati, permettendo di realizzare le operazioni richieste.
\end{definition}

Uno stesso ADT può essere realizzato mediante strutture dati
differenti, con costi temporali e spaziali diversi. Per esempio,
una coda di richieste deve elaborare gli arrivi nell'ordine in cui
si presentano: questo vincolo ne definisce il comportamento, ma
non stabilisce come memorizzare le richieste.

Nella progettazione si individua prima l'ADT adatto al problema,
quindi se ne sceglie un'implementazione in base alle operazioni
più frequenti e ai relativi costi.

\section{Array e accesso diretto}
\label{sec:fi-strutture-array}

\begin{definition}[Array]
\label{def:fi-array}
Un \emph{array} è una struttura dati che memorizza elementi di
dimensione fissa in posizioni contigue di memoria, consentendo
l'accesso diretto a ciascun elemento mediante il suo indice.
\end{definition}

Considerando un array di \(n\) elementi, ciascuno di dimensione
\(s\) byte, e indicando con \(b\) l'indirizzo del primo elemento,
l'indirizzo dell'elemento di indice \(i\), per \(0\le i<n\), è
\[
  \operatorname{addr}(v[i])=b+i\,s.
\]
Il calcolo richiede un numero costante di operazioni: accesso e
aggiornamento per indice hanno costo \(\Theta(1)\), nel modello
della \cref{sec:fi-complessita-modello}. Un array di \(n\) celle
occupa \(ns\) byte per gli elementi, ai quali si aggiungono le
informazioni di gestione.

L'organizzazione contigua favorisce la località spaziale descritta
nella \cref{sec:fi-memoria-archiviazione}. La \emph{lista} di Python
della \cref{sec:fi-python-liste} è realizzata, in CPython, mediante
un array dinamico di riferimenti agli oggetti: sono contigui i
riferimenti, non necessariamente gli oggetti stessi.

\subsection{Array statici e dinamici}
\label{sec:fi-array-dinamici}

Un array \emph{statico} ha dimensione fissata, mentre un array
\emph{dinamico} può adattare la capacità di memorizzazione al numero
di elementi presenti. Per quest'ultimo si distinguono la
\emph{dimensione} \(n\), cioè il numero di elementi presenti,
e la \emph{capacità} \(C\ge n\), cioè il numero di celle allocate.
Se ciascuna cella occupa \(s\) byte, l'area allocata richiede
\(Cs\) byte, anche quando \(n<C\).

Quando \(n=C\), un inserimento può richiedere un'area più grande,
la copia degli \(n\) elementi e il rilascio della precedente:
il singolo ridimensionamento costa \(\Theta(n)\).

Se la capacità viene aumentata di un fattore costante maggiore di
uno, il costo dei ridimensionamenti può essere distribuito fra
gli inserimenti. Nel caso di raddoppio, il costo complessivo delle
copie su \(m\ge1\) inserimenti in coda è limitato dalla somma
geometrica
\[
  1+2+4+\cdots+2^{\lfloor\log_2 m\rfloor}<2m.
\]
Considerando anche gli inserimenti stessi, il costo totale delle
\(m\) operazioni è quindi \(\Theta(m)\), e il costo
\emph{ammortizzato} di ciascuna è \(\Theta(1)\), anche se una
singola operazione che provoca il ridimensionamento può costare
\(\Theta(n)\). Il costo ammortizzato si riferisce a una sequenza
di operazioni e non presuppone una distribuzione probabilistica
degli ingressi.

Inserire un elemento all'indice \(i\) richiede di spostare verso
destra i \(n-i\) elementi successivi; rimuovere quello di indice
\(i\) richiede di spostare a sinistra i \(n-i-1\) elementi seguenti.
A capacità sufficiente, i costi sono rispettivamente
\(\Theta(n-i+1)\) e \(\Theta(n-i)\). L'inserimento in coda
è un caso particolare, al quale si applica il costo ammortizzato
precedente.

\section{Liste concatenate}
\label{sec:fi-liste-concatenate}

\begin{definition}[Lista concatenata]
\label{def:fi-lista-concatenata}
Una \emph{lista concatenata} (\emph{linked list}) è una struttura
dati composta da nodi, ciascuno contenente un elemento e uno o
più riferimenti ad altri nodi. I nodi non devono occupare
posizioni contigue di memoria.
\end{definition}

Nella \emph{lista semplicemente concatenata}, ogni nodo contiene
un riferimento \texttt{next} al successivo. Un riferimento
\texttt{head} individua il primo nodo, mentre l'ultimo ha
\texttt{next} nullo. Se la lista è vuota, anche \texttt{head}
è nullo.

\begin{figure}[htbp]
  \centering
  \begin{tikzpicture}[unipd/diagram,x=1cm,y=1cm]
  \node[anchor=east,font=\small] at (-4.4,0.8) {Array};
  \foreach \x/\v in {-3.6/A,-2.15/B,-0.7/C,0.75/D} {
    \node[unipd/node,minimum width=1cm,minimum height=.65cm] at (\x,0.8) {\v};
  }
  \node[anchor=east,font=\small] at (-4.4,-0.9) {Lista concatenata};
  \foreach \x/\v/\id in {-3.1/A/a,-1.2/B/b,0.7/C/c,2.6/D/d} {
    \node[unipd/node,minimum width=1.25cm,minimum height=.65cm] (\id) at (\x,-0.9) {\v\,\(\mid\)\,\scriptsize next};
  }
  \node[font=\scriptsize] (head) at (-3.1,-1.65) {\texttt{head}};
  \draw[unipd/directed edge] (head.north) -- (a.south);
  \draw[unipd/directed edge] (a) -- (b);
  \draw[unipd/directed edge] (b) -- (c);
  \draw[unipd/directed edge] (c) -- (d);
  \node[font=\scriptsize] (null) at (4.35,-0.9) {\texttt{None}};
  \draw[unipd/directed edge] (d) -- (null);
\end{tikzpicture}

  \caption{Organizzazione contigua di un array e collegamenti tra
  i nodi di una lista semplicemente concatenata.}
  \label{fig:fi-array-lista-concatenata}
\end{figure}

Nella \cref{fig:fi-array-lista-concatenata} l'array consente
l'accesso diretto mediante indice, mentre nella lista concatenata
occorre seguire i riferimenti a partire da \texttt{head}.
Per raggiungere il nodo di indice \(i\) sono necessari \(i\)
passaggi: considerando anche l'accesso al nodo iniziale, il costo
è \(\Theta(i+1)\), e nel caso peggiore \(\Theta(n)\).
L'attraversamento dei nodi mediante i riferimenti successivi è
detto anche \emph{link hopping}.

\subsection{Inserimento, rimozione e varianti}
\label{sec:fi-liste-concatenate-operazioni}

Inserire un nodo all'inizio richiede di collegarlo alla vecchia
testa e aggiornare \texttt{head}. Per rimuovere la testa,
quando la lista non è vuota, è sufficiente spostare \texttt{head}
al nodo successivo. Entrambe le operazioni hanno costo
\(\Theta(1)\), anche quando l'inserimento avviene in una lista
inizialmente vuota.

Se è noto il riferimento a un nodo \(p\), un nuovo nodo \(q\)
può essere inserito \emph{dopo} \(p\) mediante i due aggiornamenti
\[
  q.\texttt{next}\gets p.\texttt{next},
  \qquad p.\texttt{next}\gets q.
\]
Analogamente, per rimuovere il nodo \emph{successivo} a \(p\),
se esiste, basta porre
\[
  p.\texttt{next}\gets p.\texttt{next.next}.
\]
L'aggiornamento dei collegamenti richiede \(\Theta(1)\) tempo;
se il nodo di riferimento deve prima essere cercato, si aggiunge
il costo dell'attraversamento. In particolare, conoscere il nodo
da rimuovere non è in generale sufficiente per eliminarlo in
tempo costante da una lista semplicemente concatenata: può
essere necessario individuare il suo predecessore.

Un riferimento aggiuntivo \texttt{tail} all'ultimo nodo rende
costante l'inserimento in fondo, purché venga aggiornato quando
necessario. Non consente invece, da solo, di rimuovere l'ultimo
nodo in tempo costante, poiché occorre individuare il predecessore.

In una \emph{lista doppiamente concatenata} ciascun nodo conserva
anche il riferimento \texttt{prev} al precedente, permettendo
l'attraversamento in entrambe le direzioni. Se il nodo da
eliminare è già individuato, i collegamenti ai nodi adiacenti
possono essere aggiornati in \(\Theta(1)\) tempo; rimane invece
necessario il costo della ricerca quando il nodo non è noto.
In una \emph{lista circolare}, il collegamento dell'ultimo nodo
riporta al primo. Queste varianti modificano le operazioni
disponibili e la memoria necessaria per i riferimenti.

\begin{table}[htbp]
  \centering
  \small
  \begin{tabular}{lll}
    \hline
    Operazione & Array dinamico & Lista concatenata\\
    \hline
    Accesso all'indice \(i\) & \(\Theta(1)\) & \(\Theta(i+1)\)\\
    Inserimento in testa & \(\Theta(n+1)\) & \(\Theta(1)\)\\
    Rimozione in testa & \(\Theta(n)\) & \(\Theta(1)\)\\
    Inserimento in posizione nota & \(O(n)\) & \(\Theta(1)\)\\
    Inserimento in fondo & \(O(1)\) ammort. & \(\Theta(1)\) con \texttt{tail}\\
    \hline
  \end{tabular}
  \caption{Costi delle principali operazioni su array dinamici
  e liste semplicemente concatenate, con \(n\) elementi
  e indici da zero.}
  \label{tab:fi-array-concatenata-costi}
\end{table}

Nella riga relativa a una posizione nota, l'array riceve
l'indice di inserimento, mentre per la lista concatenata
è già disponibile il riferimento al nodo precedente.
Le liste concatenate sono adatte a inserimenti e rimozioni
presso nodi già individuati, mentre gli array sono preferibili
quando prevalgono gli accessi per indice. La scelta deve
considerare anche la memoria aggiuntiva per i riferimenti
e la località spaziale della \cref{sec:fi-memoria-archiviazione}.

\section{Pila: politica LIFO}
\label{sec:fi-adt-pila}

\begin{definition}[Pila]
\label{def:fi-adt-pila}
Una \emph{pila} (\emph{stack}) è un ADT in cui inserimento,
rimozione e ispezione interessano soltanto la cima, secondo la
politica \emph{LIFO} (\emph{Last In, First Out}).
\end{definition}

Le operazioni fondamentali sono \texttt{push(e)}, che inserisce
l'elemento \(e\) in cima, \texttt{pop()}, che rimuove e restituisce
la cima, e \texttt{top()}, che restituisce la cima senza rimuoverla.
Le operazioni \texttt{is\_empty()} e \texttt{size()} verificano
rispettivamente se la pila è vuota e quanti elementi contiene.
Le operazioni \texttt{pop()} e \texttt{top()} richiedono una pila
non vuota: il comportamento in caso contrario deve essere
specificato dall'interfaccia.

\begin{figure}[htbp]
  \begin{minipage}[c]{0.54\textwidth}
    La \cref{fig:fi-pila} mostra una pila contenente tre elementi.
    L'unica estremità accessibile è la cima, sulla quale agiscono
    \texttt{push}, \texttt{pop} e \texttt{top}. L'ultimo elemento
    inserito è quindi il primo che può essere rimosso, secondo la
    politica LIFO.
  \end{minipage}
  \hfill
  \begin{minipage}[c]{0.40\textwidth}
    \centering
    \begin{tikzpicture}[unipd/diagram,x=1cm,y=1cm]
  \foreach \y/\v/\id in {0/10/a,0.75/20/b,1.5/30/c} {
    \node[unipd/node,minimum width=.85cm,
      minimum height=.65cm] (\id) at (0,\y) {\v};
  }

  \node[font=\scriptsize] at (0,2.05) {cima};
  \node[font=\scriptsize] at (0,-.55) {base};

  \draw[unipd/directed edge]
    (c.east) -- ++(1.3,0)
    node[midway,above=3pt,font=\scriptsize]
      {\texttt{push(e)}}
    node[midway,below=3pt,font=\scriptsize]
      {\texttt{pop()}};
\end{tikzpicture}

  \end{minipage}
  \caption{Rappresentazione schematica di una pila.}
  \label{fig:fi-pila}
\end{figure}

La pila delle chiamate della \cref{sec:fi-pila-ricorsione} segue
questa stessa politica. Altre applicazioni comprendono la gestione
dell'annullamento delle operazioni e il ritorno a posizioni
precedenti durante una navigazione.

\subsection{Implementazione della pila mediante array}
\label{sec:fi-adt-pila-implementazione-array}

Con un array dinamico, la cima coincide con l'ultimo elemento
presente. Si può conservarne esplicitamente l'indice
\texttt{topIdx}, inizialmente \(-1\) nella pila vuota e aggiornato
a ogni inserimento o rimozione. L'operazione \texttt{push}
aggiunge in fondo, mentre \texttt{pop} rimuove l'ultimo elemento:
entrambe richiedono \(O(1)\) tempo ammortizzato.
Le operazioni \texttt{top()}, \texttt{is\_empty()} e
\texttt{size()} richiedono \(\Theta(1)\).

Una realizzazione in Python utilizza direttamente queste operazioni
sulla lista:

\begin{unipdcode}[language=Python]
class Pila:
    def __init__(self):
        self._dati = []

    def push(self, elemento):
        self._dati.append(elemento)

    def pop(self):
        if self.is_empty():
            raise IndexError("Pila vuota")
        return self._dati.pop()

    def top(self):
        if self.is_empty():
            raise IndexError("Pila vuota")
        return self._dati[-1]

    def is_empty(self):
        return len(self._dati) == 0

    def size(self):
        return len(self._dati)
\end{unipdcode}


Un'implementazione alternativa utilizza una lista semplicemente
concatenata, nella quale la cima coincide con \texttt{head}:
inserimento, rimozione e lettura in testa richiedono
\(\Theta(1)\). Se si mantiene un contatore aggiornato a ogni
modifica, anche \texttt{size()} costa \(\Theta(1)\); senza
contatore, richiede \(\Theta(n)\) per attraversare la lista.

\section{Coda: politica FIFO}
\label{sec:fi-adt-coda}

\begin{definition}[Coda]
\label{def:fi-adt-coda}
Una \emph{coda} (\emph{queue}) è un ADT in cui gli inserimenti
avvengono in fondo e le rimozioni all'inizio, secondo la politica
\emph{FIFO} (\emph{First In, First Out}).
\end{definition}

Le operazioni fondamentali sono \texttt{enqueue(e)}, che inserisce
l'elemento \(e\) in fondo, \texttt{dequeue()}, che rimuove e
restituisce il primo elemento, e \texttt{first()}, che lo
restituisce senza rimuoverlo. Le operazioni \texttt{is\_empty()}
e \texttt{size()} verificano rispettivamente se la coda è vuota
e quanti elementi contiene. Le operazioni \texttt{dequeue()} e
\texttt{first()} richiedono una coda non vuota: il comportamento
in caso contrario deve essere specificato dall'interfaccia.

\begin{figure}[htbp]
  \begin{minipage}[c]{0.54\textwidth}
    La \cref{fig:fi-coda} mostra una coda contenente tre elementi.
    Gli inserimenti avvengono in fondo e le rimozioni all'inizio:
    il primo elemento inserito è quindi il primo che può essere
    rimosso, secondo la politica FIFO.
  \end{minipage}
  \hfill
  \begin{minipage}[c]{0.40\textwidth}
    \centering
    \resizebox{\linewidth}{!}{%
      \begin{tikzpicture}[unipd/diagram,x=1cm,y=1cm]
  \foreach \x/\v/\id in {-1.15/A/a,0/B/b,1.15/C/c} {
    \node[unipd/node,minimum width=.85cm,
      minimum height=.7cm] (\id) at (\x,0) {\v};
  }
  \node[font=\scriptsize] at (-1.15,-.65) {inizio};
  \node[font=\scriptsize] at (1.15,-.65) {fondo};
  \draw[unipd/directed edge]
    (-1.63,0) -- (-2.4,0);
  \node[font=\scriptsize] at (-2.15,.55)
    {\texttt{dequeue()}};
  \draw[unipd/directed edge]
    (2.4,0) -- (1.63,0);
  \node[font=\scriptsize] at (2.15,.55)
    {\texttt{enqueue(e)}};
\end{tikzpicture}
%
    }
  \end{minipage}
  \caption{Rappresentazione schematica di una coda.}
  \label{fig:fi-coda}
\end{figure}

Le code di stampa e l'elaborazione dei pacchetti di rete sono
applicazioni della politica FIFO.

\subsection{Implementazione della coda mediante array circolare}
\label{sec:fi-coda-circolare}

Un array dinamico ordinario con rimozione dal primo indice
richiederebbe di spostare tutti gli elementi successivi, con
costo \(\Theta(n)\). Per evitarlo si utilizza un
\emph{array circolare} con capacità \(C>0\), indice \(f\) del
primo elemento e numero \(n\) di elementi presenti.
La posizione logica \(i\) occupa l'indice fisico
\[
  (f+i)\bmod C,
  \qquad 0\le i<n.
\]
In particolare, se \(n<C\), l'indice del prossimo inserimento è
\[
  b=(f+n)\bmod C.
\]
L'inserimento scrive in \(b\) e incrementa \(n\); la rimozione
legge l'elemento in \(f\), libera la posizione occupata, pone
\(f\gets(f+1)\bmod C\) e decrementa \(n\). L'operatore modulo
permette di riutilizzare le posizioni liberate senza spostare
gli elementi, come illustra la \cref{fig:fi-coda-circolare}.

\begin{figure}[htbp]
  \centering
  \begin{tikzpicture}[unipd/diagram,x=1cm,y=1cm]
  \foreach \i/\v in {0/E,1/F,2/{},3/{},4/C,5/D} {
    \node[unipd/node,minimum width=.85cm,minimum height=.68cm] (c\i) at ({(\i-2.5)*1.15},0) {\v};
    \node[font=\scriptsize] at ({(\i-2.5)*1.15},-0.65) {\i};
  }
  \node[font=\small] (front) at (1.725,1.25) {\texttt{front}=4};
  \draw[unipd/directed edge] (front) -- (c4.north);
  \node[font=\small] (back) at (-.575,1.25) {\texttt{back}=2};
  \draw[unipd/directed edge] (back) -- (c2.north);
  \node[font=\small] at (0,-1.22) {Ordine logico: \(C,D,E,F\)};
\end{tikzpicture}

  \caption{Coda circolare di capacità sei: \(f=4\), \(n=4\)
  e posizione del prossimo inserimento \(b=2\).}
  \label{fig:fi-coda-circolare}
\end{figure}

Quando la capacità è esaurita, il ridimensionamento ricopia
gli elementi nell'ordine logico a partire dal primo e aggiorna
gli indici. Con la strategia della \cref{sec:fi-array-dinamici},
\texttt{enqueue()} richiede \(O(1)\) tempo ammortizzato, mentre
\texttt{dequeue()}, \texttt{first()}, \texttt{is\_empty()}
e \texttt{size()} richiedono \(\Theta(1)\).

La seguente implementazione in Python mantiene l'indice del
primo elemento e il numero di elementi presenti. Quando la
capacità è esaurita, raddoppia l'array ricopiando gli elementi
nell'ordine logico.

\begin{unipdcode}[language=Python]
class CodaCircolare:
    def __init__(self):
        self._dati = [None] * 4
        self._inizio = 0
        self._dimensione = 0

    def is_empty(self):
        return self._dimensione == 0

    def size(self):
        return self._dimensione

    def enqueue(self, elemento):
        capacita = len(self._dati)
        if self._dimensione == capacita:
            nuovo = [None] * (2 * capacita)
            for i in range(self._dimensione):
                nuovo[i] = self._dati[(self._inizio + i) % capacita]
            self._dati = nuovo
            self._inizio = 0
            capacita *= 2
        indice = (self._inizio + self._dimensione) % capacita
        self._dati[indice] = elemento
        self._dimensione += 1

    def first(self):
        if self.is_empty():
            raise IndexError("Coda vuota")
        return self._dati[self._inizio]

    def dequeue(self):
        elemento = self.first()
        self._dati[self._inizio] = None
        self._inizio = (self._inizio + 1) % len(self._dati)
        self._dimensione -= 1
        return elemento
\end{unipdcode}


\subsection{Implementazione della coda mediante lista concatenata}
\label{sec:fi-coda-concatenata}

Una lista semplicemente concatenata può mantenere sia il
riferimento \texttt{head}, corrispondente all'inizio, sia
\texttt{tail}, corrispondente all'ultimo nodo.
L'operazione \texttt{enqueue()} collega il nuovo nodo dopo
\texttt{tail}, mentre \texttt{dequeue()} rimuove la testa.
Entrambe richiedono \(\Theta(1)\) tempo, purché i riferimenti
vengano aggiornati correttamente anche durante l'inserimento
del primo elemento e la rimozione dell'ultimo.

Mantenendo un contatore aggiornato a ogni modifica, anche
\texttt{size()} richiede \(\Theta(1)\); senza contatore,
richiede \(\Theta(n)\) per attraversare la lista. In assenza
del riferimento \texttt{tail}, l'inserimento in fondo richiede
invece \(\Theta(n)\).

\section{Coda doppia}
\label{sec:fi-adt-deque}

\begin{definition}[Coda doppia]
\label{def:fi-adt-deque}
Una \emph{coda doppia} (\emph{deque}, da \emph{double-ended queue})
è un ADT che consente inserimenti, rimozioni e ispezioni a entrambe
le estremità, senza prevedere l'accesso arbitrario agli elementi
interni.
\end{definition}

Le operazioni fondamentali sono \texttt{add\_first(e)} e
\texttt{add\_last(e)}, che inseriscono l'elemento \(e\)
rispettivamente all'inizio e alla fine, e
\texttt{delete\_first()} e \texttt{delete\_last()},
che rimuovono e restituiscono gli elementi alle stesse estremità.
Le operazioni \texttt{first()} e \texttt{last()} restituiscono
i rispettivi elementi senza rimuoverli, mentre
\texttt{is\_empty()} e \texttt{size()} verificano se la coda
è vuota e quanti elementi contiene. Le operazioni di rimozione
e ispezione richiedono una coda non vuota: il comportamento
in caso contrario deve essere specificato dall'interfaccia.

La coda doppia generalizza la pila e la coda delle
\cref{sec:fi-adt-pila,sec:fi-adt-coda}. Limitando le operazioni
a una sola estremità si ottiene la politica LIFO, mentre
inserendo in fondo e rimuovendo all'inizio si ottiene
la politica FIFO.

Una coda doppia può essere realizzata mediante un array circolare
ridimensionabile oppure una lista doppiamente concatenata con
riferimenti a entrambe le estremità. Nel primo caso inserimenti
e rimozioni richiedono \(O(1)\) tempo ammortizzato; nel secondo
richiedono \(\Theta(1)\). Con entrambe le rappresentazioni,
le operazioni di ispezione e verifica dello stato richiedono
\(\Theta(1)\), mantenendo aggiornato il numero di elementi.

Fra le applicazioni figurano la gestione di attività da entrambe
le estremità di una sequenza e la verifica dei palindromi,
confrontando e rimuovendo progressivamente il primo e
l'ultimo elemento.

\section{Lista come tipo di dato astratto}
\label{sec:fi-adt-lista}

\begin{definition}[Lista]
\label{def:fi-adt-lista}
Una \emph{lista} è un ADT che rappresenta una sequenza di elementi
in cui la posizione di ciascuno è significativa. Gli elementi
possono essere letti, aggiornati, inseriti e rimossi mediante
indici, e possono comparire più volte nella stessa lista.
\end{definition}

Le operazioni fondamentali sono \texttt{get(i)}, che restituisce
l'elemento di indice \(i\), \texttt{set(i,e)}, che lo sostituisce
con \(e\), e \texttt{add(i,e)}, che inserisce \(e\) nella posizione
indicata, spostando gli elementi successivi. L'operazione
\texttt{append(e)} aggiunge un elemento in fondo, mentre
\texttt{remove(i)} rimuove e restituisce quello di indice \(i\),
spostando a sinistra gli elementi successivi.
Le operazioni \texttt{is\_empty()} e \texttt{size()} verificano
rispettivamente se la lista è vuota e quanti elementi contiene.
Per una lista di \(n\) elementi, gli indici validi per lettura,
aggiornamento e rimozione appartengono a \([0,n-1]\), mentre
per l'inserimento è ammesso anche \(i=n\).

L'ADT lista può essere realizzato mediante un array dinamico
oppure una lista concatenata, introdotti nelle
\cref{sec:fi-strutture-array,sec:fi-liste-concatenate}.
La scelta dipende dalle operazioni più frequenti, come mostra
la \cref{tab:fi-array-concatenata-costi}: l'array consente
l'accesso diretto mediante indice, mentre una lista concatenata
permette inserimenti e rimozioni in tempo costante quando
sono già disponibili i riferimenti necessari.

La lista Python della \cref{sec:fi-python-liste} realizza
queste funzionalità mediante un array dinamico, ma presenta
un'interfaccia differente: \texttt{insert(i, e)} inserisce
un elemento nella posizione indicata, \texttt{pop(i)}
rimuove e restituisce quello di indice \(i\), mentre
\texttt{remove(e)} elimina la prima occorrenza del valore \(e\).
Le operazioni e i relativi costi dipendono quindi dalla
struttura dati adottata, non dalla sola specifica astratta.


\section{Mappa}
\label{sec:fi-adt-mappa}

\begin{definition}[Mappa]
\label{def:fi-adt-mappa}
Una \emph{mappa} (\emph{map}) è un ADT che rappresenta una
collezione di associazioni tra chiavi distinte e valori.
Ogni chiave identifica un solo valore, mentre chiavi diverse
possono essere associate allo stesso valore.
\end{definition}

Le operazioni fondamentali sono \texttt{put(k,v)}, che inserisce
una nuova associazione o aggiorna il valore di una chiave esistente,
\texttt{get(k)}, che restituisce il valore associato alla chiave,
e \texttt{remove(k)}, che elimina l'associazione.
L'operazione \texttt{contains(k)} verifica la presenza di una chiave,
mentre \texttt{keys()}, \texttt{values()} e \texttt{items()}
consentono di accedere rispettivamente alle chiavi, ai valori
e alle associazioni presenti. Le operazioni \texttt{is\_empty()}
e \texttt{size()} verificano se la mappa è vuota e quante
associazioni contiene. Per \texttt{get(k)} e \texttt{remove(k)},
il comportamento in caso di chiave assente deve essere
specificato dall'interfaccia.

\begin{example}[Conteggio delle frequenze]
Per ogni parola \(w\) letta in un testo, se \(w\) non è presente
nella mappa si registra l'associazione \(w\mapsto1\); altrimenti
si incrementa di uno il valore associato alla chiave.
Il risultato associa a ciascuna parola il numero delle sue
occorrenze.
\end{example}

L'ADT mappa non richiede che le chiavi siano ordinate.
Le complessità delle operazioni dipendono dalla struttura dati
utilizzata per realizzarlo, poiché l'interfaccia astratta non
prescrive una particolare implementazione. In Python, il tipo
\texttt{dict} fornisce un'implementazione di questo comportamento.

\section{Insieme}
\label{sec:fi-adt-insieme}

\begin{definition}[Insieme]
\label{def:fi-adt-insieme}
Un \emph{insieme} (\emph{set}) è un ADT che rappresenta una
collezione di elementi distinti, senza associare loro valori
né prevedere un ordinamento.
\end{definition}

Le operazioni fondamentali sono \texttt{add(e)}, che inserisce
l'elemento \(e\) se assente, \texttt{remove(e)}, che lo elimina,
e \texttt{contains(e)}, che ne verifica l'appartenenza.
Il comportamento di \texttt{remove(e)} quando l'elemento è assente
deve essere specificato dall'interfaccia.

Si considerano inoltre le operazioni di unione, intersezione
e differenza:
\[
  A\cup B,\qquad A\cap B,\qquad A\setminus B,
\]
che restituiscono rispettivamente gli elementi presenti in almeno
uno dei due insiemi, in entrambi oppure soltanto nel primo.
La verifica di inclusione \(A\subseteq B\) può essere esposta
mediante \texttt{is\_subset(T)}, che controlla se tutti gli
elementi dell'insieme appartengono anche a \texttt{T}.
Completano l'interfaccia la consultazione degli elementi e le
operazioni \texttt{is\_empty()} e \texttt{size()}, che verificano
se l'insieme è vuoto e quanti elementi contiene.
L'ordine di enumerazione non fa parte della specifica astratta.

\begin{example}[Eliminazione dei duplicati]
Inserendo successivamente i valori \(4,2,2,8,6,2\) in un insieme,
si ottiene \(\{2,4,6,8\}\). Una successiva inserzione di \(2\)
non ne modifica il contenuto. Su insiemi di documenti, unione,
intersezione e differenza permettono inoltre di combinare risultati
di ricerca in base alla presenza di determinati termini.
\end{example}

L'ADT insieme si distingue dalla mappa della
\cref{sec:fi-adt-mappa}: nel primo gli elementi sono unici,
nella seconda lo sono le chiavi, ma non necessariamente i valori.
In Python, il tipo \texttt{set} fornisce un'implementazione
di questo comportamento. Le complessità delle operazioni
dipendono dalla struttura dati adottata, non dalla specifica
astratta dell'ADT.
